\slcol{श्रीभगवानुवाच —\\
\Index{इमं विवस्वते योगं} प्रोक्तवानहमव्ययम् ।\\
विवस्वान्मनवे प्राह मनुरिक्ष्वाकवेऽब्रवीत् ॥ ४.१ ॥}
\translate{śrībhagavānuvāca — \\
imaṁ vivasvatē yōgaṁ prōktavānahamavyayam | \\
vivasvānmanavē prāha manurikṣvākavē'bravīt ‖ 4.1 ‖} 
\cquote{Sri Bhagavan said-\\
This imperishable \emph{yoga} was imparted by Me to Vivasvan, the Sun God. Vivasvan conveyed it to Manu, and Manu passed it on to Ikshvaku.}

\slcol{\Index{एवं परम्पराप्राप्तमिमं} राजर्षयो विदुः । \\
स कालेनेह महता योगो नष्टः परन्तप ॥ ४.२ ॥ }
\translate{ēvaṁ paramparāprāptamimaṁ rājarṣayō viduḥ | \\
sa kālēnēha mahatā yōgō naṣṭaḥ parantapa ‖ 4.2 ‖ }
\cquote{This ancient knowledge was handed down in an unbroken lineage, and the great sages knew it well. But over time, O Parantapa (scorcher of foes), this \emph{yoga} has been lost in the world.}



\newpage
\begin{mananam}{\mananamfont {Mananam sloka - 1, 2}} 
\mananamtext What is my understanding of the ancient tradition? How can I better appreciate the preservation of such teachings handed on through a succession of \emph{rishis} and teachers who have committed their lives to this? What are the benefits I see in my life, or can hope to see, by embracing these teachings? 
\end{mananam}
\WritingHand\enspace{\selfReflect\small{Self Reflection}}
\begin{inspiration}{\mananamfont Inspiration}
\mananamtext In modern times, where society needs scientific proof for everything, we might question the historic reality of many of the characters in our ancient texts. The proof of the Gita is not to be sought in history or science. Instead, the proof of the Gita is to be found by applying its teachings in our life and by examining whether it gives us mental peace and clarity. To focus our attention on petty outer details is to miss the inner transformational message that has been handed down to us.
\end{inspiration}
\newpage

\slcol{\Index{स एवायं मया} तेऽद्य योगः प्रोक्तः पुरातनः ।\\
भक्तोऽसि मे सखा चेति रहस्यं ह्येतदुत्तमम् ॥ ४.३ ॥}
\translate{sa ēvāyaṁ mayā tē'dya yōgaḥ prōktaḥ purātanaḥ |\\
bhaktō'si mē sakhā cēti rahasyaṁ hyētaduttamam ‖ 4.3 ‖}
\cquote{That ancient and supreme \emph{yoga} has today been declared to you by Me, for you are My devotee and My friend—thus, it is a profound secret indeed.}
\slcol{अर्जुन उवाच —\\
\Index{अपरं भवतो जन्म} परं जन्म विवस्वतः ।\\
 कथमेतद्विजानीयां त्वमादौ प्रोक्तवानिति ॥ ४.४ ॥}
\translate{Arjuna uvāca —\\
aparaṁ bhavatō janma paraṁ janma vivasvataḥ |\\
kathamētadvijānīyāṁ tvamādau prōktavāniti ‖ 4 .4 ‖}
\cquote{Arjuna said-\\
Your birth was in recent times, while Vivasvan’s was much prior to that. How then am I to understand that You taught this \emph{yoga} at the beginning?}

\slcol{श्रीभगवानुवाच —\\
\Index{बहूनि मे व्यतीतानि} जन्मानि तव चार्जुन । \\
तान्यहं वेद सर्वाणि न त्वं वेत्थ परन्तप ॥ ४.५ ॥ }
\translate{śrībhagavānuvāca —\\
bahūni mē vyatītāni janmāni tava cārjuna |\\
tānyahaṁ vēda sarvāṇi na tvaṁ vēttha parantapa ‖ 4.5 ‖}
\cquote{Sri Bhagavan said - You and I have passed through many births, O Arjuna. I know them all, but you do not recall them, O Parantapa (scorcher of foes).}

\slcol{\Index{अजोऽपि सन्नव्ययात्मा} भूतानामीश्वरोऽपि सन् ।\\
प्रकृतिं स्वामधिष्ठाय सम्भवाम्यात्ममायया ॥ ४.६ ॥}
\translate{ajō'pi sannavyayātmā bhūtānāmīśvarō'pi san |\\
prakṛtiṁ svāmadhiṣṭhāya sambhavāmyātmamāyayā ‖ 4.6 ‖}
\cquote{Although I am unborn and imperishable and I am the Lord of all beings; yet, being the basis for My \emph{prakriti}, I am born of My own \emph{maya}.}


\newpage
\begin{mananam}{\mananamfont {Mananam sloka - 5}} 
\mananamtext What is my understanding of the theory of reincarnation? What are my doubts about it? Can the understanding of reincarnation and the law of \emph{karma} be useful in directing my actions in this life? Can this concept help me to see an underlying balance and justice in society?
\end{mananam}
\WritingHand\enspace{\selfReflect\small{Self Reflection}}
\begin{inspiration}{\mananamfont Inspiration}
\mananamtext Reincarnation is not hard to believe when we see child prodigies in any field. How did their talents and skills come about? 
The fact that our memory is erased in every life helps us to start fresh in this life to change our negative tendencies. If the past is the blueprint for the present, then the present offers an opportunity for us to shape a better future. 
\end{inspiration}
\newpage

\slcol{\Index{यदा यदा हि धर्मस्य} ग्लानिर्भवति भारत ।\\
अभ्युत्थानमधर्मस्य तदात्मानं सृजाम्यहम् ॥ ४.७ ॥ }
\translate{yadā yadā hi dharmasya glānirbhavati bhārata | \\
abhyutthānamadharmasya tadātmānaṁ sṛjāmyaham ‖ 4.7 ‖ }
\cquote{Whenever there is a decline in \emph{dharma} (virtue) and a rise in \emph{adharma} (vice), O Bhārata (Arjuna), then I manifest Myself.
}
\slcol{\Index{परित्राणाय साधूनां} विनाशाय च दुष्कृताम् ।\\
धर्मसंस्थापनार्थाय सम्भवामि युगे युगे ॥ ४.८ ॥}
\translate{paritrāṇāya sādhūnāṁ vināśāya ca duṣkṛtām |\\
dharmasaṁsthāpanārthāya sambhavāmi yugē yugē ‖ 4.8 ‖}
\cquote{For the protection of the good and holy, for the destruction of the wicked, and for the re-establishment of dharma (righteousness), I take birth in every era (\emph{yuga})}

\slcol{\Index{जन्म कर्म च मे} दिव्यमेवं यो वेत्ति तत्त्वतः ।\\
 त्यक्त्वा देहं पुनर्जन्म नैति मामेति सोऽर्जुन ॥ ४.९ ॥}
\translate{janma karma ca mē divyamēvaṁ yō vētti tattvataḥ |\\
tyaktvā dēhaṁ punarjanma naiti māmēti sō'rjuna ‖ 4.9 ‖}
\cquote{One who truly understands the divine nature of My birth and actions does not take birth again after leaving the body; rather, he attains Me, O Arjuna.}

\slcol{\Index{वीतरागभयक्रोधा} मन्मया मामुपाश्रिताः ।\\
बहवो ज्ञानतपसा पूता मद्भावमागताः ॥ ४.१० ॥ }
\translate{vītarāgabhayakrōdhā manmayā māmupāśritāḥ |\\
bahavō jñānatapasā pūtā madbhāvamāgatāḥ ‖ 4.10 ‖}
\cquote{Freed from attachment, fear, and anger; absorbed in Me and taking refuge in Me; purified by the austerities of knowledge—many have attained to My being.}



\newpage
\begin{mananam}{\mananamfont {Mananam sloka - 7, 8}} 
\mananamtext What is my response to injustices in society? How do I feel when I see or encounter an individual or a group being unfair and seemingly getting away with impunity? Can I find solace in the Lord’s promise? Can I use it to find a sense of inner acceptance when faced with circumstances wherein outer change is not possible?
\end{mananam}
\WritingHand\enspace{\selfReflect\small{Self Reflection}}
\begin{inspiration}{\mananamfont Inspiration}
\mananamtext Our concept of a higher power in which we believe must be one that is not just loving and kind but also powerful and just. Then such a God can be an ideal source of refuge, one that fills us with hope and strength.
\end{inspiration}
\newpage


\newpage
\begin{mananam}{\mananamfont {Mananam sloka - 10}} 
\mananamtext Does attaining the Lord symbolise being in the Divine state? Can I relate to such a mental state wherein fear, anger and such emotions do not trouble me?
What does it mean to be purified by the fire of knowledge? How can knowledge help me when strong emotions bring out the dark side of my personality?
\end{mananam}
\WritingHand\enspace{\selfReflect\small{Self Reflection}}
\begin{inspiration}{\mananamfont Inspiration}
\mananamtext When the dormant afflictions (\emph{keshas}) become active in us, we become somewhat possessed by a different personality. Our ability to function with our higher wisdom becomes dormant, thus making us somewhat unconscious of our own natural state. Awareness is the first light of knowledge that guides us back home, to our true nature. 
\end{inspiration}
\newpage

\slcol{\Index{ये यथा मां प्रपद्यन्ते} तांस्तथैव भजाम्यहम् ।\\
मम वर्त्मानुवर्तन्ते मनुष्याः पार्थ सर्वशः ॥ ४.११ ॥}
\translate{yē yathā māṁ prapadyantē tāṁstathaiva bhajāmyaham |\\
mama vartmānuvartantē manuṣyāḥ pārtha sarvaśaḥ ‖ 4.11 ‖}
\cquote{In whatever way people approached Me, so do I respond to them. O Pārtha (Arjuna), all human beings follow My path in every way.}
\slcol{\Index{काङ्क्षन्तः कर्मणां} सिद्धिं यजन्त इह देवताः ।\\ 
क्षिप्रं हि मानुषे लोके सिद्धिर्भवति कर्मजा ॥ ४.१२ ॥}
\translate{kāṅkṣantaḥ karmaṇāṁ siddhiṁ yajanta iha dēvatāḥ |\\
kṣipraṁ hi mānuṣē lōkē siddhirbhavati karmajā ‖ 4.12 ‖}
\cquote{Those who desire success in this world offer sacrifices to the gods; for human beings attain fulfilment quickly through action.}

\slcol{\Index{चातुर्वर्ण्यं मया सृष्टं} गुणकर्मविभागशः ।\\
 तस्य कर्तारमपि मां विद्ध्यकर्तारमव्ययम् ॥ ४.१३ ॥}
\translate{cāturvarṇyaṁ mayā sṛṣṭaṁ guṇakarmavibhāgaśaḥ |\\
tasya kartāramapi māṁ viddhyakartāramavyayam ‖ 4.13 ‖}
\cquote{The four-fold social order was created by Me, based on the division of \emph{guṇa} (qualities) and \emph{karma} (duties). Though I am the agent (of this division), know Me to be the non-doer and immutable.}

\slcol{\Index{न मां कर्माणि लिम्पन्ति} न मे कर्मफले स्पृहा ।\\
इति मां योऽभिजानाति कर्मभिर्न स बध्यते ॥ ४.१४ ॥}
\translate{na māṁ karmāṇi limpanti na mē karmaphalē spṛhā |\\
iti māṁ yō'bhijānāti karmabhirna sa badhyatē ‖ 4.14 ‖}
\cquote{I am not bound by actions, nor do I desire their fruits. One who knows Me thus is not bound by actions.}



\newpage
\begin{mananam}{\mananamfont {Mananam sloka - 11}} 
\mananamtext What is my concept of God or the Supreme entity? Am I able to have faith in a higher force in the universe or in a purer self within me? Can I learn from observation that everyone in the world has their own faith in a concept of God? Can I also observe in my own life as to how any concept faith is developed, sustained and strengthened?
\end{mananam}
\WritingHand\enspace{\selfReflect\small{Self Reflection}}
\begin{inspiration}{\mananamfont Inspiration}
\mananamtext It might appear paradoxical that the Lord proclaims that He strengthens the belief of everyone by rewarding them, regardless of what concept of God they hold onto. The God of \emph{Sanatana Dharma} is not a jealous God but empowers everyone to follow their own faith. As followers of this eternal truth, we should be open-minded and tolerant towards others of different beliefs. 
\end{inspiration}
\newpage




\newpage
\begin{mananam}{\mananamfont {Mananam sloka - 14}} 
\mananamtext Do I realise the importance of doing my duties without leaving any trace of the imprint of those actions in my mind? What does it mean to me to be free from the effects of any actions I do? 
Do I harbour a desire for personal gratification from my actions? What is the trap of seeking personal gratification?
\end{mananam}
\WritingHand\enspace{\selfReflect\small{Self Reflection}}
\begin{inspiration}{\mananamfont Inspiration}
\mananamtext The Lord has shared with us a formula for freedom, the one which He Himself uses: When one becomes free from the desire for results, there is no imprint or \emph{karmic} blueprint of any action, nor of action’s consequences, in one’s own mind.
\end{inspiration}
\newpage

\slcol{\Index{एवं ज्ञात्वा कृतं} कर्म पूर्वैरपि मुमुक्षुभिः ।\\
कुरु कर्मैव तस्मात्त्वं पूर्वैः पूर्वतरं कृतम् ॥ ४.१५ ॥}
\translate{ēvaṁ jñātvā kṛtaṁ karma pūrvairapi mumukṣubhiḥ |\\
kuru karmaiva tasmāttvaṁ pūrvaiḥ pūrvataraṁ kṛtam ‖ 4.15 ‖}
\cquote{Having known this, the seekers of liberation in ancient times also performed action. Therefore, you too should act, just as the ancients did in former times.}
\slcol{\Index{किं कर्म किमकर्मेति} कवयोऽप्यत्र मोहिताः ।\\
तत्ते कर्म प्रवक्ष्यामि  यज्ज्ञात्वा मोक्ष्यसेऽशुभात् ॥ ४.१६ ॥}
\translate{kiṁ karma kimakarmēti  kavayō'pyatra mōhitāḥ |\\
tattē karma pravakṣyāmi  yajjñātvā mōkṣyasē'śubhāt ‖ 4.16 ‖}
\cquote{Even the wise are deluded about what constitutes action and inaction. Therefore, I shall explain to you the nature of action (the truth of action and inaction); understanding this, you shall be freed from evil (of \emph{saṁsāra}).}

\slcol{\Index{कर्मणो ह्यपि बोद्धव्यं} बोद्धव्यं च विकर्मणः ।\\
अकर्मणश्च बोद्धव्यं गहना कर्मणो गतिः ॥ ४.१७ ॥}
\translate{karmaṇō hyapi bōddhavyaṁ bōddhavyaṁ ca vikarmaṇaḥ |\\
akarmaṇaśca bōddhavyaṁ gahanā karmaṇō gatiḥ ‖ 4.17 ‖}
\cquote{One must understand the nature of action, as well as that of forbidden action (\emph{vikarma}) and inaction (\emph{akarma}); for the true nature of action is indeed difficult to comprehend.}

\slcol{\Index{कर्मण्यकर्म यः} पश्येदकर्मणि च कर्म यः ।\\
स बुद्धिमान्मनुष्येषु स युक्तः कृत्स्नकर्मकृत् ॥ ४.१८ ॥}
\translate{karmaṇyakarma yaḥ paśyēdakarmaṇi ca karma yaḥ |\\
sa buddhimānmanuṣyēṣu sa yuktaḥ kṛtsnakarmakṛt ‖ 4.18 ‖}
\cquote{He who sees inaction in action, and action in inaction—he is wise among men; he is a \emph{yogi} and the true performer of all actions.}



\newpage
\begin{mananam}{\mananamfont {Mananam sloka - 16, 17}} 
\mananamtext How do I differentiate between right action and wrong action? What are my parameters? Are they based only on what is good for me and my family, or are they based on universal values? How can I shift my modus operandi to working based on spiritual values? Am I tuned in to my conscience when I have to make difficult decisions in my life?
\end{mananam}
\WritingHand\enspace{\selfReflect\small{Self Reflection}}
\begin{inspiration}{\mananamfont Inspiration}
\mananamtext Within this very life, the results of wrong actions can be seen as creating addictions, stress, negativity, and conflicts. Results of right actions are mentally uplifting in this life as well as future lives, if any. 
\end{inspiration}
\newpage


\newpage
\begin{mananam}{\mananamfont {Mananam sloka - 18}} 
\mananamtext Whenever I get confused between right and wrong actions, do I then tend to sit back, fearing to take any action? Do I choose inaction because I am lazy or afraid of consequences? Do I introspect into what causes my resistance to my duties, self-improvement, and spiritual practices?
Have I experienced what it is to work effortlessly? If so, what was it like and what were the factors that led to such a state? Can I relate to the state wherein actions that are free of self are empowering and uplifting?
\end{mananam}
\WritingHand\enspace{\selfReflect\small{Self Reflection}}
\begin{inspiration}{\mananamfont Inspiration}
\mananamtext The secret to effortless action is to immerse oneself entirely in one’s duties and to work from that state. Where there is no self, there the energy of the Divine flows in full force. To not resist the flow of the Divine is the best way to live life.
\end{inspiration}
\newpage

\slcol{\Index{यस्य सर्वे समारम्भाः} कामसङ्कल्पवर्जिताः ।\\
ज्ञानाग्निदग्धकर्माणं तमाहुः पण्डितं बुधाः ॥ ४.१९ ॥}
\translate{yasya sarvē samārambhāḥ kāmasaṅkalpavarjitāḥ |\\
jñānāgnidagdhakarmāṇaṁ tamāhuḥ paṇḍitaṁ budhāḥ ‖ 4.19 ‖}
\cquote{One whose actions are entirely free from desire and personal motive, and whose deeds have been consumed by the fire of knowledge—such a person is declared by the wise to be a sage.}
\slcol{\Index{त्यक्त्वा कर्मफलासङ्गं} नित्यतृप्तो निराश्रयः ।\\
कर्मण्यभिप्रवृत्तोऽपि नैव किञ्चित्करोति सः ॥ ४.२० ॥}
\translate{tyaktvā karmaphalāsaṅgaṁ nityatṛptō nirāśrayaḥ |\\
karmaṇyabhipravṛttō'pi naiva kiñcitkarōti saḥ ‖ 4.20 ‖}
\cquote{Having renounced attachment to the fruits of action, ever content and dependent on nothing—such a person, though fully engaged in action, truly does nothing at all.}
 
\slcol{\Index{निराशीर्यतचित्तात्मा} त्यक्तसर्वपरिग्रहः ।\\
शारीरं केवलं कर्म कुर्वन्नाप्नोति किल्बिषम् ॥ ४.२१ ॥}
\translate{nirāśīryatacittātmā tyaktasarvaparigrahaḥ |\\
śārīraṁ kēvalaṁ karma kurvannāpnōti kilbiṣam ‖ 4.21 ‖}
\cquote{Free from desires, with mind and self under control, having renounced all attachment to possessions—such a one acts merely with the body and incurs no sin.}

\slcol{\Index{यदृच्छालाभसन्तुष्टो} द्वन्द्वातीतो विमत्सरः ।\\ 
समः सिद्धावसिद्धौ च कृत्वापि न निबध्यते ॥ ४.२२ ॥}
\translate{yadṛcchālābhasantuṣṭō dvandvātītō vimatsaraḥ |\\
samaḥ siddhāvasiddhau ca kṛtvāpi na nibadhyatē ‖ 4.22 ‖}
\cquote{Content with whatever comes by chance, transcending the dualities, free from envy, and even-minded in success and failure—such a one, though engaged in action, is not bound by it.}

\newpage
\begin{mananam}{\mananamfont {Mananam sloka - 19}} 
\mananamtext What is the relation between wisdom and the actions of my daily life? When I reflect on the past, can I observe that the nature of action at any point in my life was based on the degree of wisdom I then had? How can I operate from a higher level of wisdom to guide my actions? 
Does being free from personal desires improve the mental and spiritual quality of all my works?
\end{mananam}
\WritingHand\enspace{\selfReflect\small{Self Reflection}}
\begin{inspiration}{\mananamfont Inspiration}
\mananamtext When we are obsessed with our desires, we get stressed out and block our own path. When we truly want something harmless and meaningful to our lives, the best course is to make a wish or prayer and then mentally relax, trusting in God and the forces of universe to bring unto us what is right for us.
\end{inspiration}
\newpage


\newpage
\begin{mananam}{\mananamfont {Mananam sloka - 22}} 
\mananamtext How do I respond to success and failure? Does success make me elated and failure disappoint me? Am I content to accept what comes naturally in my life rather than mentally hankering after what others have? Do I evaluate myself based on others’ successes? Is it possible to be content in life while yet being active with one’s duties?
\end{mananam}
\WritingHand\enspace{\selfReflect\small{Self Reflection}}
\begin{inspiration}{\mananamfont Inspiration}
\mananamtext A sign of spiritual understanding is to not evaluate oneself based on one’s material success and achievements. Being content in what one gets in one’s life and to fulfil all of one’s duties without comparing oneself to others is real freedom. Spiritual wisdom is that which brings out such a state of inner freedom.
\end{inspiration}
\newpage

\slcol{\Index{गतसङ्गस्य मुक्तस्य} ज्ञानावस्थितचेतसः ।\\
यज्ञायाचरतः कर्म समग्रं प्रविलीयते ॥ ४.२३ ॥}
\translate{gatasaṅgasya muktasya jñānāvasthitacētasaḥ |\\
yajñāyācarataḥ karma samagraṁ pravilīyatē ‖ 4.23 ‖}
\cquote{Free from attachment, the liberated one whose mind is established in knowledge, and who performs action as \emph{yajña} (offering to the Lord)—all their actions (with their results) are dissolved.}
\slcol{\Index{ब्रह्मार्पणं ब्रह्म} हविर्ब्रह्माग्नौ ब्रह्मणा हुतम् ।\\
ब्रह्मैव तेन गन्तव्यं ब्रह्मकर्मसमाधिना ॥ ४.२४ ॥}
\translate{brahmārpaṇaṁ brahma havirbrahmāgnau brahmaṇā hutam |\\
brahmaiva tēna gantavyaṁ brahmakarmasamādhinā ‖ 4.24 ‖}
\cquote{The oblation is Brahman (the absolute and supreme reality), the ghee offered is Brahman, the sacrificial fire is Brahman, and the one who makes the offering is also Brahman. One who perceives Brahman in every action attains only Brahman.}

\slcol{\Index{दैवमेवापरे यज्ञं} योगिनः पर्युपासते ।\\
ब्रह्माग्नावपरे यज्ञं यज्ञेनैवोपजुह्वति ॥ ४.२५ ॥}
\translate{daivamēvāparē yajñaṁ yōginaḥ paryupāsatē |\\
brahmāgnāvaparē yajñaṁ yajñēnaivōpajuhvati ‖ 4.25 ‖}
\cquote{Some \emph{yogis} offer sacrifice to the \emph{devas} (gods); others offer the self (ego-identification) as a sacrifice by the self (pure Consciousness), into the fire of Brahman alone.}

\slcol{\Index{श्रोत्रादीनीन्द्रियाण्यन्ये} संयमाग्निषु जुह्वति ।\\
शब्दादीन्विषयानन्य इन्द्रियाग्निषु जुह्वति ॥ ४.२६ ॥}
\translate{śrōtrādīnīndriyāṇyanyē saṁyamāgniṣu juhvati |\\
śabdādīnviṣayānanya indriyāgniṣu juhvati ‖ 4.26 ‖}
\cquote{Some again offer the organ of hearing and other senses as sacrifice into the fire of restraint (symbolising mastery of the mind over the senses); others offer sound and other sense objects as sacrifice into the fire of the senses (symbolising controlled engagement with sense objects).}

\newpage
\begin{mananam}{\mananamfont {Mananam sloka - 24}} 
\mananamtext How can my life be an offering to the Divine? Can all of my thought, speech, and action be imbued with a sense of offering? Can I see the value in reminding myself every day to do all actions in this spirit of sacrificing the little ego-self? What are the ways I can apply such an internal practice in my daily life? 
\end{mananam}
\WritingHand\enspace{\selfReflect\small{Self Reflection}}
\begin{inspiration}{\mananamfont Inspiration}
\mananamtext Offering all of one’s actions to the Supreme Spirit and seeing everything as non-distinct from that One Consciousness is the most exalted mind state. It frees one from all stresses in life and immediately rewards one with mental peace. 
\end{inspiration}
\newpage

\slcol{\Index{सर्वाणीन्द्रियकर्माणि} प्राणकर्माणि चापरे ।\\
आत्मसंयमयोगाग्नौ जुह्वति ज्ञानदीपिते ॥ ४.२७ ॥}
\translate{sarvāṇīndriyakarmāṇi prāṇakarmāṇi cāparē |\\
ātmasaṁyamayōgāgnau juhvati jñānadīpitē ‖ 4.27 ‖}
\cquote{Others again, who have mastered all sensory activities and the functions of \emph{prāṇa} (vital energies), offer them into the fire of the \emph{yoga} of self-restraint, kindled by knowledge.}
\slcol{\Index{द्रव्ययज्ञास्तपोयज्ञा} योगयज्ञास्तथापरे ।\\
स्वाध्यायज्ञानयज्ञाश्च यतयः संशितव्रताः ॥ ४.२८ ॥}
\translate{dravyayajñāstapōyajñā yōgayajñāstathāparē |\\
svādhyāyajñānayajñāśca yatayaḥ saṁśitavratāḥ ‖ 4.28 ‖}
\cquote{Some offer wealth, austerities, and \emph{yoga} as sacrifice; while the ascetics of self-restraint and rigid vows offer self-study of the scriptures and knowledge as sacrifice.}

\slcol{\Index{अपाने जुह्वति प्राणं} प्राणेऽपानं तथापरे ।\\
प्राणापानगती रुद्ध्वा प्राणायामपरायणाः ॥ ४.२९ ॥}
\translate{apānē juhvati prāṇaṁ prāṇē'pānaṁ tathāparē |\\
prāṇāpānagatī ruddhvā prāṇāyāmaparāyaṇāḥ ‖ 4.29 ‖}
\cquote{Some offer the outgoing breath (\emph{prāṇa}) into the incoming breath (\emph{apāna}), and the incoming breath into the outgoing breath; regulating both, they remain devoted to breath control (\emph{prāṇāyāma}).
}

\slcol{\Index{अपरे नियताहाराः} प्राणान्प्राणेषु जुह्वति ।\\
सर्वेऽप्येते यज्ञविदो यज्ञक्षपितकल्मषाः ॥ ४.३० ॥}
\translate{aparē niyatāhārāḥ prāṇānprāṇēṣu juhvati |\\
sarvē'pyētē yajñavidō yajñakṣapitakalmaṣāḥ ‖ 4.30 ‖}
\cquote{Others, who regulate their diet, offer the vital breaths (\emph{prāṇa}) into the controlled vital breaths. All of them are knowers of sacrifice, and through sacrifice, their sins are destroyed.}
\newpage

\slcol{\Index{यज्ञशिष्टामृतभुजो यान्ति} ब्रह्म सनातनम् ।\\
नायं लोकोऽस्त्ययज्ञस्य कुतोऽन्यः कुरुसत्तम ॥ ४.३१ ॥}
\translate{yajñaśiṣṭāmṛtabhujō yānti brahma sanātanam |\\
nāyaṁ lōkō'styayajñasya kutō'nyaḥ kurusattama ‖ 4.31 ‖}
\cquote{Those who partake of the nectar-like remnants of sacrifice attain the eternal Brahman. But those who do not perform sacrifice do not gain even this world—what then of the other (higher worlds), O best among the Kurus (Arjuna)?}

\slcol{\Index{एवं बहुविधा यज्ञा} वितता ब्रह्मणो मुखे ।\\
कर्मजान्विद्धि तान्सर्वानेवं ज्ञात्वा विमोक्ष्यसे ॥ ४.३२ ॥}
\translate{ēvaṁ bahuvidhā yajñā vitatā brahmaṇō mukhē |\\
karmajānviddhi tānsarvānēvaṁ jñātvā vimōkṣyasē ‖ 4.32 ‖} 
\cquote{Thus, various kinds of sacrifices are spread out at the mouth of the vedas (are known through the vedas). Know them all to be born of action; and knowing this, you shall attain liberation.}

\slcol{\Index{श्रेयान्द्रव्यमयाद्यज्ञा}ज्ज्ञानयज्ञः परन्तप ।\\
सर्वं कर्माखिलं पार्थ ज्ञाने परिसमाप्यते ॥ ४.३३ ॥}
\translate{śrēyāndravyamayādyajñājjñānayajñaḥ parantapa |\\
sarvaṁ karmākhilaṁ pārtha jñānē parisamāpyatē ‖ 4.33 ‖}
\cquote{The sacrifice of knowledge is superior to the sacrifice of material offerings, O Parantapa (scorcher of enemies). All actions, in their entirety, O Pārtha, culminate in knowledge.}

\slcol{\Index{तद्विद्धि प्रणिपातेन} परिप्रश्नेन सेवया ।\\
उपदेक्ष्यन्ति ते ज्ञानं ज्ञानिनस्तत्त्वदर्शिनः ॥ ४.३४ ॥}
\translate{tadviddhi praṇipātēna paripraśnēna sēvayā |\\
upadēkṣyanti tē jñānaṁ jñāninastattvadarśinaḥ ‖ 4.34 ‖}
\cquote{Know this through wholehearted prostration, sincere inquiry, and devoted service to the wise. Those who have realised the truth shall impart that knowledge to you.}



\newpage
\begin{mananam}{\mananamfont {Mananam sloka - 32, 33}} 
\mananamtext What are the various ways I seek to appease worldly authorities to gain personal favour? Similarly, do I worship gods and deities for the attainment of material benefits in life? The scriptures and sages say that it is superior to seek the Supreme Spirit and to be established in the Divine state than to seek after hundreds of worldly goals which are impermanent in nature. What do I need to do to embrace such values and perspective in my life? How can I prioritise seeking goals that are long-term, inwardly uplifting, and beneficial for others as well?
\end{mananam}
\WritingHand\enspace{\selfReflect\small{Self Reflection}}
\begin{inspiration}{\mananamfont Inspiration}
\mananamtext The only real value of sacrifices, worship and other rituals are to lift us out of excessive obsession to worldly pursuits. With the purity of mind gained through surrender to higher spiritual forces, we learn to value the pursuit of wisdom and work towards realizing our true self.  
\end{inspiration}
\newpage


\newpage
\begin{mananam}{\mananamfont {Mananam sloka - 34}} 
\mananamtext Do I approach teachers for seeking wisdom and spiritual truth? Do I have resistance to approaching them with humility? Am I willing to pose my questions in a reverential manner? Am I generous enough to offer something of use to them? Can I make time to be of service to them or to humanity, thus affirming my spiritual life as my priority?
\end{mananam}
\WritingHand\enspace{\selfReflect\small{Self Reflection}}
\begin{inspiration}{\mananamfont Inspiration}
\mananamtext A seeker of spiritual truth must first develop certain skills in appropriately approaching the teacher and of being service to him. These are not just cultural or social norms, but a mental attitude that makes the transmission of knowledge possible.  
\end{inspiration}
\newpage

\slcol{\Index{यज्ज्ञात्वा न पुनर्मोहमेवं} यास्यसि पाण्डव ।\\
येन भूतान्यशेषेण द्रक्ष्यस्यात्मन्यथो मयि ॥ ४.३५ ॥}
\translate{yajjñātvā na punarmōhamēvaṁ yāsyasi pāṇḍava |\\
yēna bhūtānyaśēṣēṇa drakṣyasyātmanyathō mayi ‖ 4.35 ‖}
\cquote{O Pandava (Arjuna), by knowing this, you will not come under delusion like this again. Through this knowledge, you will behold all beings, without exception, in the self and in Me.}
\slcol{\Index{अपि चेदसि पापेभ्यः} सर्वेभ्यः पापकृत्तमः ।\\
सर्वं ज्ञानप्लवेनैव वृजिनं सन्तरिष्यसि ॥ ४.३६ ॥}
\translate{api cēdasi pāpēbhyaḥ sarvēbhyaḥ pāpakṛttamaḥ |\\
sarvaṁ jñānaplavēnaiva vṛjinaṁ santariṣyasi ‖ 4.36 ‖}
\cquote{Even if you were to be the most sinful among all sinners, yet you shall cross over all the (sea of) sins through the raft of knowledge (for such is its power)}

\slcol{\Index{यथैधांसि समिद्धो}ऽग्निर्भस्मसात्कुरुतेऽर्जुन ।\\
ज्ञानाग्निः सर्वकर्माणि भस्मसात्कुरुते तथा ॥ ४.३७ ॥}
\translate{yathaidhāṁsi samiddhō'gnirbhasmasātkurutē'rjuna |\\
jñānāgniḥ sarvakarmāṇi bhasmasātkurutē tathā ‖ 4.37 ‖}
\cquote{Just as a blazing fire reduces firewood to ashes, O Arjuna, so too does the fire of knowledge reduce all actions(and their binding results) into ashes.}

\slcol{\Index{न हि ज्ञानेन सदृशं} पवित्रमिह विद्यते ।\\
तत्स्वयं योगसंसिद्धः कालेनात्मनि विन्दति ॥ ४.३८ ॥}
\translate{na hi jñānēna sadṛśaṁ pavitramiha vidyatē |\\
tatsvayaṁ yōgasaṁsiddhaḥ kālēnātmani vindati ‖ 4.38 ‖}
\cquote{Truly, there is nothing here as purifying as knowledge. One perfected in \emph{yoga} realises That, in time, within one’s own self.}


\newpage
\begin{mananam}{\mananamfont {Mananam sloka - 36, 37}} 
\mananamtext Do I have hidden feelings of remorse and guilt for certain actions I did in the past? Do I tend to suppress those feelings? Can there be a healthy means of addressing them? How can I use spiritual wisdom to address such feelings in my life? Have I found a means to apply the spiritual wisdom in my life in such a way that I break away from old negative patterns of thoughts and actions, and embrace new positive actions?
\end{mananam}
\WritingHand\enspace{\selfReflect\small{Self Reflection}}
\begin{inspiration}{\mananamfont Inspiration}
\mananamtext When we realize our true nature and stop yielding to our lower animalistic impulses, we uplift the overall mental and spiritual quality of our life. Then the past can no longer haunt us or pull us down.
\end{inspiration}
\newpage


\newpage
\begin{mananam}{\mananamfont {Mananam sloka - 38}} 
\mananamtext Am I aware that just as I seek knowledge through books, teachers and other mediums, so also I can seek knowledge within myself? What is my understanding of this process of seeking knowledge inwardly versus seeking it outwardly? Does it mean to just know about my body and its various processes? Or is it about knowing the mind and various mental states? Or could it involve something much more profound which I am yet to discover?
\end{mananam}
\WritingHand\enspace{\selfReflect\small{Self Reflection}}
\begin{inspiration}{\mananamfont Inspiration}
\mananamtext Self-knowledge is the greatest purifier. All spiritual practices are a means to bring us to this. Until one is established in this wisdom, one should not prematurely give up one’s duties and practices in life.    
\end{inspiration}
\newpage

\slcol{\Index{श्रद्धावांल्लभते ज्ञानं} तत्परः संयतेन्द्रियः ।\\
ज्ञानं लब्ध्वा परां शान्तिमचिरेणाधिगच्छति ॥ ४.३९ ॥}
\translate{śraddhāvāṁllabhatē jñānaṁ tatparaḥ saṁyatēndriyaḥ |\\
jñānaṁ labdhvā parāṁ śāntimacirēṇādhigacchati ‖ 4.39 ‖}
\cquote{The one who is full of faith, devoted to knowledge, and has mastered the senses, gains knowledge. Upon attaining it, one immediately realises the supreme peace.}
\slcol{\Index{अज्ञश्चाश्रद्दधानश्च} संशयात्मा विनश्यति ।\\
नायं लोकोऽस्ति न परो न सुखं संशयात्मनः ॥ ४.४० ॥}
\translate{ajñaścāśraddadhānaśca saṁśayātmā vinaśyati |\\
nāyaṁ lōkō'sti na parō na sukhaṁ saṁśayātmanaḥ ‖ 4.40 ‖}
\cquote{The ignorant, the faithless, and the doubting self is doomed. For such a sceptic, there is neither this world, nor the next, nor any happiness.}

\slcol{\Index{योगसंन्यस्तकर्माणं} ज्ञानसञ्छिन्नसंशयम् ।\\
आत्मवन्तं न कर्माणि निबध्नन्ति धनञ्जय ॥ ४.४१ ॥}
\translate{yōgasaṁnyastakarmāṇaṁ jñānasañchinnasaṁśayam |\\
ātmavantaṁ na karmāṇi nibadhnanti dhanañjaya ‖ 4.41 ‖}
\cquote{The one who has renounced actions (and their fruits) through \emph{yoga}, whose doubts have been dispelled by knowledge, and who abides in the self (vigilantly), O Dhananjaya (Arjuna), is not bound by action.}

\slcol{\Index{तस्मादज्ञानसम्भूतं} हृत्स्थं ज्ञानासिनात्मनः ।\\
छित्त्वैनं संशयं योगमातिष्ठोत्तिष्ठ भारत ॥ ४.४२ ॥}
\translate{tasmādajñānasambhūtaṁ hṛtsthaṁ jñānāsinātmanaḥ |\\
chittvainaṁ saṁśayaṁ yōgamātiṣṭhōttiṣṭha bhārata ‖ 4.42 ‖}
\cquote{Therefore, with the sword of self-knowledge, cut asunder the doubts born of ignorance that dwell in your heart. Take refuge in \emph{yoga}, O Arjuna—and arise!}

\newpage
\begin{mananam}{\mananamfont {Mananam sloka - 39, 40}} 
\mananamtext What is the relationship between faith and knowledge? How can doubts and lack of self-confidence be destructive? How should questions and doubts be resolved before they become destructive and harmful? 
How can knowledge bring me to a state of supreme peace? What is the nature of such knowledge which requires to subdue the senses? Can meditation be a means to train the senses and lead me to peace?
\end{mananam}
\WritingHand\enspace{\selfReflect\small{Self Reflection}}
\begin{inspiration}{\mananamfont Inspiration}
\mananamtext The greatest disease of humankind is to doubt one’s own self. Even the gods cannot help such a one. Doubts are useful when they motivate one towards inquiry but harmful if they make one sceptical and negative. 
As long as there is lack of self-knowledge, one is bound to remain in delusion, which leads to a life of struggle and despair. 
\end{inspiration}

\newpage
\begin{center}
  {\devfont ॐ तत्सदिति श्रीमद्भगवद्गीतासूपनिषत्सु ब्रह्मविद्यायां योगशास्त्रे\\
  श्रीकृष्णार्जुनसम्वादे ज्ञानकर्मसन्न्यासयोगो नाम चतुर्थोऽध्यायः॥\\}
  \chapEndSlokaTranslate{ōṃ tatsaditi śrīmadbhagavadgītās ūpaniṣatsu \\brahmavidyāyāṃ yōgaśāstrē\\ śrīkṛṣṇārjunasaṃvādē jñāna-karma-sannyāsa-yoga\\nāma caturtho'dhyāyaḥ‖}
\end{center}

